% Dictionarium Sinicum — NAACL/COLING 2027 Findings via ARR October 2026 cycle
% Previous round: ARR March 2026 (OpenReview sVPHR4PSzR, avg 3.33 Borderline).
% This revision addresses reviewer objections; response inlined as Appendix D.
% Compile with LuaLaTeX (required for CJK characters)
% lualatex dictionarium_sinicum && bibtex dictionarium_sinicum && lualatex dictionarium_sinicum && lualatex dictionarium_sinicum
\documentclass[11pt]{article}

% Anonymous review mode for ARR submission
\usepackage[review]{acl}

% microtype with modest character protrusion — reduces bad line-breaks
% without destabilizing the layout (SamJ, ARR March 2026 feedback).
\usepackage[activate={true,nocompatibility},final,factor=1050]{microtype}
\usepackage{booktabs}  % Professional table rules
\usepackage{multirow}

% Require at least 3 characters on both sides of any hyphenation break.
% Kills all 2-letter fragments ("in-", "re-", "en-", "es-", "ex-", etc.)
% that dominated the ARR March 2026 PDF. Wrapped in \AtBeginDocument
% because babel resets these primitives on \selectlanguage.
% NOTE: this reflows body content to 8 pages (from 7). ACL/EMNLP long
% papers allow 8 body pages + unlimited refs/appendix, so this is fine.
\AtBeginDocument{%
  \lefthyphenmin=3\relax
  \righthyphenmin=3\relax
}

% \hyphenation directives — words WITHOUT a '-' get no break at all;
% words WITH '-' pin the allowed break points.
\hyphenation{%
  % Proper nouns and headers — never break these.
  Chinese Cantonese Spanish Portuguese Vietnamese Tibetan European
  German French Arabic Japanese Korean Table Section Entries Creative
  Conference Association Language Languages SambaNova BabelNet DeepSeek
  Hispanico Community Sinicum Kirsten Ponzetto Association
  % Common words that always broke ugly — never break.
  published entry entries dialect dialects system context format
  corpus primary program numbered normalizer sentences terminology
  discarded complete compilation consolidated contamination
  correct disabled single dictionary prevent formatting
  lexicon lexical specialized inference specific
  covering coverage generates generate
  MiniMax Commercial commercial combined currently filter genuine
  maritime backfill produces produced producing
  % Where breaks are unavoidable, pin them at natural syllable boundaries.
  lex-i-cog-ra-phy
  magni-tude
  defi-ni-tions
  defi-ni-tion
  compa-ra-ble
  multi-lin-gual
  medi-cal
}

% Font setup for LuaLaTeX — supports CJK + Arabic + Cyrillic
\usepackage[english,bidi=default]{babel}
\babelfont{rm}{TeXGyreTermesX}  % Times-compatible

% CJK support
\usepackage{luatexja-fontspec}
\setmainjfont{Noto Serif CJK SC}  % Chinese (Simplified)

% Prevent luatexja from inserting inter-character spacing around Cyrillic
\ltjsetparameter{jacharrange={-2}}

% Korean
\babelprovide[import]{korean}
\babelfont[*hangul]{rm}{Noto Serif CJK KR}

% Arabic script (for Persian examples)
\babelprovide[import]{arabic}
\babelfont[*arabic]{rm}{Noto Sans Arabic}

% Thai script (for loanword examples)
\babelprovide[import]{thai}
\babelfont[*thai]{rm}{Noto Sans Thai}

% Devanagari (for Hindi examples)
\babelprovide[import]{hindi}
\babelfont[*hindi]{rm}{Noto Sans Devanagari}

% Verbatim environments for prompts
\usepackage{fancyvrb}
\usepackage{listings}
\lstset{
  basicstyle=\ttfamily\small,
  breaklines=true,
  frame=single,
  xleftmargin=0.5em,
  xrightmargin=0.5em,
}

% For the title
\setlength\titlebox{6.5cm}

\title{Dictionarium Sinicum: Building a 428K-Headword Chinese Dictionary\\in 18 Languages with a Single LLM for Under \$150}

\author{Anonymous}

\begin{document}
\maketitle

\begin{abstract}
We present \textit{Dictionarium Sinicum}, an open multilingual Chinese dictionary comprising 428,073 headwords with definitions in 18 languages, constructed by merging seven community-maintained dictionaries with glosses generated by a single open-weight LLM (MiniMax M2.5) at a total API inference cost of \${\raise.17ex\hbox{$\scriptstyle\sim$}}146 at published pay-as-you-go rates. A three-phase pipeline---batch translation, retry, and context-aware backfill---achieves 100\% coverage across all 7.7M headword--language slots. We describe prompt engineering techniques that prevent script contamination in CJK-target languages, including Japanese kanji echo and Korean hanja mixing, and a deterministic post-translation script validator that detects and rejects cross-language contamination at a rate of 0.83\% of all generated definitions. Against four community reference dictionaries the pipeline reaches 87.3\% sense coverage; when compared to ten MT baselines given a generic prompt, the pipeline leads by 12--22 percentage points on this metric, though a substantial fraction of that gap reflects the pipeline's access to community-reference glosses in its input context (\S\ref{sec:community}). Even the \textit{same model} with a generic zero-context prompt scores 14 points lower than our pipeline, isolating a real prompt-engineering effect. The dictionary, augmented with 511K Cantonese and Hokkien dialect forms, is released under CC~BY-SA~4.0.
\end{abstract}

%% ============================================================
\section{Introduction}
\label{sec:intro}

Bilingual and multilingual dictionaries remain essential infrastructure for language learning, translation, and NLP\@. For Chinese, the open-source community has produced high-quality dictionaries for a handful of language pairs---CC-CEDICT for English ({\raise.17ex\hbox{$\scriptstyle\sim$}}124K entries), HanDeDict for German ({\raise.17ex\hbox{$\scriptstyle\sim$}}159K entries), CFDICT for French ({\raise.17ex\hbox{$\scriptstyle\sim$}}56K entries), and CC-CIDICT for Indonesian ({\raise.17ex\hbox{$\scriptstyle\sim$}}124K entries)---but no comparable resources exist for Korean, Russian, Vietnamese, Tagalog, Persian, Swedish, Dutch, Portuguese, Arabic, Thai, Hindi, or Italian (CHEDICC provides a small Spanish dictionary of {\raise.17ex\hbox{$\scriptstyle\sim$}}4,400 entries). Creating such dictionaries manually requires years of expert effort per language pair.

Recent advances in large language models suggest an alternative: using a single multilingual model to generate dictionary-quality definitions at scale. However, prior work on LLM-assisted lexicography has been limited to small-scale experiments. The closest precedent is a Pre-Qin philosophy lexicon project using Qwen3-14B to generate bilingual dictionary entries for approximately 1,000 classical Chinese philosophical terms in a single language pair \citep{wu2025preqin}. \citet{lu2024chainofdict} use existing bilingual dictionaries as chain-of-thought context to improve LLM translation of rare words---a complementary approach where dictionaries help LLMs translate, whereas we use LLMs to \textit{build} dictionaries. No published work has attempted dictionary generation at the scale of hundreds of thousands of headwords across more than two target languages.

We address this gap with \textit{Dictionarium Sinicum}, a system that:

\begin{enumerate}
\item \textbf{Merges} seven community-maintained Chinese dictionaries (CC-CEDICT, HanDeDict, CFDICT, CC-CIDICT, CHEDICC, Wiktextract, JMdict) covering 6 language pairs into a unified headword table of 428,073 entries.
\item \textbf{Translates} all headwords into 18 languages using a single open-weight LLM (MiniMax M2.5, available on multiple inference providers) with carefully engineered prompts that produce dictionary-style output and prevent script contamination.
\item \textbf{Backfills} gaps using a context-aware prompt that sends existing translations as read-only context, ensuring consistency across languages.
\item \textbf{Augments} the dictionary with 511,514 dialect forms from 18 open sources covering Cantonese (Jyutping, 99.9\% single-character coverage) and Hokkien (POJ/Tâi-lô) pronunciations and lexical equivalents.
\end{enumerate}

Total inference cost was \${\raise.17ex\hbox{$\scriptstyle\sim$}}146 and the full pipeline completed in under one week.


%% ============================================================
\section{Related Work}
\label{sec:related}

\subsection{Multilingual Lexical Resources}

Large-scale multilingual lexical resources have been constructed through both manual and automated methods. PanLex \citep{kamholz2014panlex} aggregates bilingual dictionaries into a graph covering {\raise.17ex\hbox{$\scriptstyle\sim$}}5,700 languages with over 1.2 billion pairwise translations, but its Chinese coverage is limited to headword-level mappings without definition glosses. BabelNet \citep{navigli2012babelnet} combines WordNet, Wikipedia, and other resources into a multilingual encyclopedic dictionary, but is not open-access for bulk redistribution and focuses on synset alignment rather than dictionary-style glosses. The Open Multilingual Wordnet effort \citep{bond2012wordnets} aggregates community-contributed wordnets with per-language licence audits; our work is complementary in that it produces definition glosses at scale rather than synset-linked lexical entries. Our work differs from these prior efforts in producing complete, human-readable definition glosses for a single source language across 18 targets, released under CC~BY-SA~4.0.

\subsection{LLM-Assisted Lexicography}

The use of LLMs for dictionary construction is an emerging area with limited published work. The closest precedent is a Pre-Qin philosophy lexicon \citep{wu2025preqin} that uses Qwen3-14B to generate {\raise.17ex\hbox{$\scriptstyle\sim$}}1,000 bilingual entries for classical Chinese philosophical terms---two orders of magnitude smaller than our work. Chain-of-Dictionary prompting \citep{lu2024chainofdict} uses existing bilingual dictionaries as chain-of-thought context to improve LLM translation of rare words---a complementary approach where dictionaries help LLMs translate, whereas we use LLMs to \textit{build} dictionaries.

\subsection{Multilingual Translation Models}

General-purpose multilingual translation models such as NLLB \citep{costajussa2022nllb} cover 200+ languages but are designed for sentence-level translation, not dictionary-style glosses. Tower \citep{alves2024tower} focuses on high-resource language pairs. \citet{wang2023doclevel} show that LLMs benefit from document-level context in translation; our pipeline exploits an analogous effect at the entry level by supplying community glosses as intra-prompt context (\S\ref{sec:prompts}). Neither NLLB nor Tower has been evaluated for or applied to dictionary construction. These models do not natively produce multi-sense gloss lists, making direct comparison nontrivial (see Limitations).

\subsection{Community Chinese Dictionaries}

The open Chinese dictionary ecosystem is built on the CC-CEDICT model: community-maintained, CEDICT-format text files released under Creative Commons licenses. CC-CEDICT \citep{cedict}, HanDeDict \citep{handedict}, CFDICT \citep{cfdict}, and CC-CIDICT \citep{cidict} together cover four European languages and Indonesian. Wiktextract \citep{wiktextract} provides multilingual glosses extracted from Wiktionary, and JMdict \citep{jmdict} covers Japanese. CHEDICC provides a small Spanish resource ({\raise.17ex\hbox{$\scriptstyle\sim$}}4,400 entries), but no comparable community dictionaries exist for Korean, Russian, Vietnamese, Tagalog, Persian, Swedish, Dutch, Portuguese, Arabic, Thai, Hindi, or Italian.


%% ============================================================
\section{Data Sources and Headword Unification}
\label{sec:data}

\subsection{Source Dictionaries}

\begin{table}[t]
\centering
\small
\begin{tabular}{llrl}
\toprule
\textbf{Source} & \textbf{Lang} & \textbf{Entries} & \textbf{License} \\
\midrule
CC-CEDICT   & en & 123,127 & CC BY-SA 4.0 \\
HanDeDict   & de & 159,809 & CC BY-SA 3.0 \\
CFDICT      & fr &  56,280 & CC BY-SA 3.0 \\
CC-CIDICT   & id & 123,157 & CC BY-SA 4.0 \\
CHEDICC     & es &   4,400 & CC BY-SA 4.0 \\
Wiktextract & multi & 140,577 & CC BY-SA 4.0 \\
JMdict      & ja & 131,427 & CC BY-SA 4.0 \\
\bottomrule
\end{tabular}
\caption{Source dictionaries. All use compatible Creative Commons licenses (BY-SA 3.0/4.0).}
\label{tab:sources}
\end{table}

Table~\ref{tab:sources} lists the seven source dictionaries. All sources use compatible Creative Commons licenses, allowing the combined work to be distributed under CC~BY-SA~4.0.

\subsection{Headword Unification}

Source dictionaries use varying conventions for pinyin romanization, traditional/simplified character pairs, and entry granularity. We unified headwords on the triple (traditional, simplified, pinyin), yielding 428,073 unique headwords. Of these, 30,546 are single-character entries covering 19,471 distinct Chinese characters (many characters have multiple pinyin readings---e.g., 参 has 12 distinct entries).

\subsection{Dialect Forms}

\begin{table*}[t]
\centering
\small
\begin{tabular}{llrlp{7cm}}
\toprule
\textbf{Source} & \textbf{Dialect} & \textbf{Forms} & \textbf{License} & \textbf{Notes} \\
\midrule
rime-cantonese     & Cantonese & 110,064 & CC BY 4.0     & Largest single source; community Jyutping readings \\
CC-Canto readings  & Cantonese &  96,278 & CC BY-SA 3.0  & Jyutping overlay for CC-CEDICT headwords (Pleco) \\
CC-Canto           & Cantonese &  29,762 & CC BY-SA 3.0  & Dedicated Cantonese dictionary \citep{cccanto} \\
Unihan kCantonese  & Cantonese &  28,800 & Unicode ToU   & Unicode standard Cantonese readings \\
WikiHan            & Cant.+Hok. & 38,794 & CC BY-SA 4.0 & Wiktionary-derived readings for both dialects \\
TaiHua             & Hokkien   &  48,154 & CC BY-SA 4.0  & Taiwanese Hokkien dictionary (Tâi-lô romanization) \\
Taibun             & Hokkien   &  39,814 & MIT           & POJ-based Hokkien corpus \\
Wiktextract        & Hokkien   &  25,606 & CC BY-SA 3.0  & Hokkien readings from Wiktionary \citep{wiktextract} \\
iTaigi             & Hokkien   &  10,572 & CC BY 4.0     & Crowdsourced Taiwanese Hokkien \citep{itaigi} \\
7 additional       & Mixed     &  83,670 & Various open  & ChhoeTaigi, early missionary sources, others \\
\bottomrule
\end{tabular}
\caption{Dialect data sources. Total: 511,514 forms (290,421 Cantonese, 221,093 Hokkien) from 18 open sources. All sources use open licenses compatible with CC~BY-SA~4.0 redistribution; three sources with NC/ND restrictions (Maryknoll, Embree, Kauiokpoo) are excluded from commercial builds.}
\label{tab:dialects}
\end{table*}

We augmented the dictionary with Cantonese and Hokkien dialect data from 18 open sources (Table~\ref{tab:dialects}). Cantonese coverage is effectively complete: 16,564 unique single-character readings covering 99.9\% of BMP CJK characters. Hokkien coverage spans 6,850 single-character readings---lower in percentage terms (37\% BMP) but representing all characters with documented Hokkien pronunciations in open data.

Hokkien romanization uses two systems: Tâi-lô (the MOE standard) and traditional POJ (Peh-ōe-jī). A distinctive feature of Hokkien is its high lexical divergence from Mandarin: 68--78\% of Hokkien forms use different characters entirely (e.g., 漂亮 \textit{piàoliang} ``beautiful'' → 媠 \textit{suí}; 東西 \textit{dōngxi} ``thing'' → 物件 \textit{mih-kiānn}).

\subsection{Chinese Text Corpus}

\begin{table}[t]
\centering
\small
\begin{tabular}{lrrl}
\toprule
\textbf{Source} & \textbf{Articles} & \textbf{Chunks} & \textbf{Type} \\
\midrule
Chinese Wikipedia & 1.4M & 14M & Encyclopedia \\
Baidu Baike        & 5.6M & 33M & Encyclopedia \\
THUCNews           & 0.8M & 11M & News \\
news2016zh         & 2.5M & 18M & News \\
NiuTrans Classical & 0.3M &  4M & Literary \\
chinese-poetry     & 0.3M &  2M & Poetry \\
ChID               & 0.6M &  7M & Idiom cloze \\
LCCC               & 22M  & 24M & Dialogue \\
\midrule
\textbf{Total}     & \textbf{34M} & \textbf{113M} & \\
\bottomrule
\end{tabular}
\caption{Chinese text corpus sources. Articles are chunked at sentence boundaries into segments of 40--80 characters.}
\label{tab:corpus}
\end{table}

To provide real-world usage context during translation, we constructed a large-scale Chinese text corpus from eight publicly available sources (Table~\ref{tab:corpus}). Articles were chunked at sentence boundaries (。！？；) into segments averaging 40--80 characters. All chunks are indexed with FTS5 using a character-level CJK tokenizer, enabling sub-15ms lookup for any Chinese term. During translation, 1--2 example sentences per headword are retrieved from this corpus to provide the LLM with attested usage context (\S\ref{sec:prompts}).


%% ============================================================
\section{Translation Pipeline}
\label{sec:pipeline}

\subsection{Model Selection}

We evaluated 20+ LLMs across 10 inference providers on a 5-headword test suite covering emergency service numbers (culturally specific), anime titles (transliteration challenges), and medical terminology. Selection criteria were: (1) correct 18-language output in the prescribed format, (2) no script contamination, (3) geographic context preservation, (4) sustainable throughput at 428K headwords, and (5) cost under \$200 at published rates.

\begin{table}[t]
\centering
\small
\begin{tabular}{llrl}
\toprule
\textbf{Model} & \textbf{Provider} & \textbf{s/ent} & \textbf{Issues} \\
\midrule
MiniMax M2.5  & MiniMax   & 1.2  & None \\
Kimi K2       & Groq      & 1.8  & None \\
MiniMax M2.1  & Fireworks & 5.4  & None \\
Nemotron 253B & NIM       & 7.0  & None \\
Sonnet 4.6    & Anthropic & 19.8 & HW echo \\
DeepSeek V3.2 & NIM       & 35.3 & HW echo \\
MiniMax M2.1  & NIM       & 167  & Think ON \\
\bottomrule
\end{tabular}
\caption{Model evaluation (5-headword test). HW echo = headword echo in output. Think ON = thinking mode silently enabled.}
\label{tab:models}
\end{table}

MiniMax M2.5 was selected for production: zero issues on the test suite, 1.2s/entry throughput, no observed rate limits with 20 parallel workers, and an estimated cost of {\raise.17ex\hbox{$\scriptstyle\sim$}}\$146 for the full run at published pay-as-you-go rates (\$0.30/M input tokens, \$1.20/M output tokens; Table~\ref{tab:models}).

A notable finding is that the same open-weight model can produce dramatically different output depending on the inference provider. MiniMax M2.1 served by NVIDIA NIM took 167s per entry with thinking mode silently enabled, while the same weights on Fireworks took 5.4s with thinking disabled. Benchmark scores (SWE-bench, MMLU) showed no predictive value for this structured multilingual task: DeepSeek V3.2 (SWE-bench 73.1\%) echoed headwords, while MiniMax M2.5 (unranked) produced flawless output. These observations align with emerging evidence that inference acceleration introduces unpredictable behavioral changes \citep{kirsten2025acceleration}.

\subsection{Prompt Engineering}
\label{sec:prompts}

The production prompt generates all 18 languages simultaneously for each batch of 20 headwords. Key design decisions:

\paragraph{System prompt (176 tokens):} Establishes the lexicographer persona, output format (\texttt{xx: def1/def2}), conciseness constraints (max 5 glosses), and language-specific rules addressing three classes of script contamination:

\begin{enumerate}
\item \textbf{Japanese kanji echo:} Without explicit instruction, LLMs frequently ``define'' a Chinese word in Japanese by repeating the same characters (e.g., 銀行 → 銀行). The prompt requires: \textit{``Provide the MEANING in Japanese, not just kanji echo or kana reading.''}
\item \textbf{Korean hanja mixing:} Korean definitions often include Chinese characters (hanja) unrecognizable to most modern Korean readers. The prompt specifies: \textit{``Write in Hangul only.''}
\item \textbf{General CJK leakage:} A blanket prohibition: \textit{``Every non-Chinese definition must contain ZERO Chinese characters.''}
\end{enumerate}

\paragraph{User prompt (batch of 20):} Each entry includes traditional/simplified characters, pinyin, POS tag, all existing community definitions as context, and 1--2 example sentences drawn from the 113M-chunk Chinese corpus (Table~\ref{tab:corpus}). The corpus examples---sourced from Chinese Wikipedia, Baidu Baike, THUCNews, and news2016zh---provide real-world usage context that helps disambiguate polysemous headwords. Example sentences are retrieved via FTS5 lookup ({<}15ms per headword), filtered to prefer 20--100 character chunks.

This constitutes \textbf{pivot translation:} for the twelve languages without community dictionaries (Swedish, Korean, Russian, Vietnamese, Tagalog, Persian, Dutch, Portuguese, Arabic, Thai, Hindi, Italian), the model effectively translates from the English and other community-language glosses rather than directly from Chinese.

\paragraph{Token economics:} Approximately 5,700 tokens per batch (500 system + 2,200 input including examples + 3,000 output). Full run: 21,403 batches. At MiniMax M2.5 published rates (\$0.30/M input, \$1.20/M output), the estimated cost is {\raise.17ex\hbox{$\scriptstyle\sim$}}\$146 without prompt caching or {\raise.17ex\hbox{$\scriptstyle\sim$}}\$118 with the platform's automatic caching (\$0.03/M for cached system prompts).

\subsection{Three-Phase Execution}

\begin{table}[t]
\centering
\small
\begin{tabular}{lrrr}
\toprule
\textbf{Phase} & \textbf{Entries reached} & \textbf{Time} & \textbf{Cost} \\
\midrule
1. Main     & 374,602 & 24h & \$130 \\
2. Retry    &  53,471 & 92m & \$12 \\
3. Backfill &  82,550 & 105m & \$4 \\
\midrule
\textbf{Union} & \textbf{428,073} & \textbf{27h} & \textbf{\$146} \\
\bottomrule
\end{tabular}
\caption{Three-phase execution summary. Costs at MiniMax M2.5 published pay-as-you-go rates. ``Entries reached'' is the count of headwords \textit{touched} in that phase (not additive across phases: a partially-completed headword can be reached again in phase~2 or 3). The 7.71M total definitions after Union are produced by the union of phases, not the sum of any Defs column.}
\label{tab:phases}
\end{table}

Table~\ref{tab:phases} summarizes the three-phase execution. \textbf{Phase~1} processed all 428,073 headwords in batches of 20 with 20 parallel workers. The API returned unparseable responses for 53,471 headwords (12.5\%), primarily due to batch-level formatting errors where the model's response was cut off or malformed. \textbf{Phase~2} retried failed headwords using the same prompt, recovering 616,276 additional definitions and leaving 5,875 headwords with partial coverage. \textbf{Phase~3} used a distinct context-aware backfill prompt that sends all existing translations as read-only context, asking the model to produce \textit{only} the missing languages. This design ensures consistency: the model sees existing glosses and produces output that matches in style and specificity. Three backfill passes closed all remaining multi-character gaps, achieving 100.0\% coverage for the 397,527 multi-character headwords.

\subsection{Parser Robustness}

Parsing structured LLM output at scale required handling several failure modes: (1)~\textbf{multi-language line collapse}---the model outputs \texttt{en: bank/de: Bank/fr: banque} on a single line; a regex normalizer splits on \texttt{/xx:} boundaries; (2)~\textbf{numbered vs.\ unnumbered responses}---the backfill prompt causes the model to drop entry numbers; the parser uses blank lines as separators; (3)~\textbf{overwrite protection}---three layers prevent backfill from corrupting existing data: query filter, parser filter, and database existence check.


%% ============================================================
\section{Quality Assessment}
\label{sec:quality}

\subsection{Script Contamination Audit}

A full-corpus audit of all 7.7M minimax definitions using a deterministic script validator revealed 63,712 definitions (0.83\%) containing scripts that do not belong to the target language.

\begin{table}[t]
\centering
\small
\begin{tabular}{lrp{3.8cm}}
\toprule
\textbf{Language} & \textbf{Rate} & \textbf{Primary Contamination} \\
\midrule
Arabic    & 3.91\% & Latin (9.7K), CJK (6.1K) \\
Persian   & 1.70\% & Latin (4.5K), CJK (2.3K) \\
Japanese  & 1.15\% & Hangul (3.1K), Cyrillic (1.8K) \\
Hindi     & 0.87\% & CJK (3.0K), Arabic (0.4K) \\
Indonesian & 0.66\% & CJK (2.3K), Cyrillic (0.5K) \\
Korean    & 0.56\% & Kana (1.4K), Cyrillic (0.7K) \\
Other 12  & {0.34--0.58\%} & CJK, Cyrillic \\
\bottomrule
\end{tabular}
\caption{Script contamination by target language (full corpus). Uniformly distributed across headword ranges.}
\label{tab:contamination}
\end{table}

Table~\ref{tab:contamination} breaks down contamination by language. We identified five distinct contamination modes:

\begin{enumerate}
\item \textbf{Cyrillic-in-non-Cyrillic} ({\raise.17ex\hbox{$\scriptstyle\sim$}}10K): The model falls back to Russian when translating into less-resourced scripts. Example: Arabic for 一世 → ``\textit{zhizn'} \foreignlanguage{arabic}{واحدة}'' (Russian mixed with Arabic).
\item \textbf{Hangul-in-Japanese / Kana-in-Korean} ({\raise.17ex\hbox{$\scriptstyle\sim$}}4.4K): Systematic confusion between CJK languages.
\item \textbf{Latin-in-Arabic/Persian} ({\raise.17ex\hbox{$\scriptstyle\sim$}}14K): English or French words injected into Arabic/Persian definitions, after exempting proper names, scientific binomials, and universal acronyms.
\item \textbf{CJK-in-non-CJK} ({\raise.17ex\hbox{$\scriptstyle\sim$}}40K): Chinese characters in non-CJK definitions, after exempting metalinguistic cross-references (``variant of 迥'').
\item \textbf{Multi-language field dumps} (rare): Multiple languages written into a single field.
\end{enumerate}

\paragraph{Cleanup methodology.} We implemented a three-stage pipeline: (1)~a \textbf{deterministic script validator} that detects all Unicode script families in each definition and compares against a per-language allowlist, with exemptions for CJK in metalinguistic reference frames (detected via compiled regex patterns covering all 18 target languages) and Latin in Arabic/Persian when constituting single letters, scientific binomials, proper names, or universal acronyms; (2)~\textbf{deletion} of contaminated minimax-sourced definitions (community entries never touched); (3)~\textbf{context-aware re-translation} via the backfill pipeline with the validator integrated as a hard reject gate---contaminated output is discarded before database write.

\paragraph{Post-cleanup status.} All 428,073 headwords have complete 18-language coverage across all sources after three backfill passes. Earlier snapshots showed residual single-character gaps concentrated in Arabic, Thai, German, and Persian---obscure radicals whose most natural definitions cite Chinese characters and were initially rejected by the script validator---but subsequent backfill iterations with alternative prompting closed those gaps as well.

\subsection{Comparison with Community Dictionaries}
\label{sec:community}

For the four languages with existing community dictionaries providing glosses (English, German, French, Indonesian), we compared LLM-generated definitions against human-curated references.\footnote{JMdict was excluded because it stores kana readings rather than Japanese-language glosses.} We sampled 200 headwords stratified by community source coverage as a frequency proxy and measured two metrics across 453 headword--language pairs:

\begin{enumerate}
\item \textbf{Sense coverage:} proportion of reference senses captured by the LLM gloss (automated lexical overlap with 50\% word-overlap threshold)
\item \textbf{False sense rate:} proportion of LLM glosses not matched by any reference sense
\end{enumerate}

\paragraph{Important caveat.} The translation prompt provides all existing community definitions as input context (\S\ref{sec:prompts}). For these four languages, the model \textit{sees} the reference glosses. This comparison therefore measures \textbf{context utilization}---how faithfully the model preserves reference senses---not independent translation quality. We find that 42\% of English, 40\% of German, 55\% of French, and 60\% of Indonesian LLM glosses are substring matches of the corresponding community definition.

\begin{table}[t]
\centering
\small
\setlength{\tabcolsep}{4pt}
\begin{tabular}{llrccrr}
\toprule
\textbf{Lang} & \textbf{Ref.} & \textbf{n} & \textbf{Cov.} & \textbf{False} & \textbf{BSc} & \textbf{COMET} \\
\midrule
en & CC-CEDICT  & 139 & 85.8 & 14.1 & .932 & .777 \\
de & HanDeDict  & 105 & 87.9 & 17.9 & .905 & .712 \\
fr & CFDICT     &  70 & 81.2 & 18.6 & .926 & .787 \\
id & CC-CIDICT  & 139 & 91.3 &  3.7 & .933 & .820 \\
\midrule
\multicolumn{2}{l}{\textbf{Overall}} & \textbf{453} & \textbf{87.3} & \textbf{12.5} & \textbf{.924} & \textbf{.774} \\
\bottomrule
\end{tabular}
\caption{LLM vs.\ community dictionary. Cov./False in \%; BSc = BERTScore F1 (XLM-R-large); COMET = \texttt{wmt22-comet-da} \citep{rei2022comet}. Semantic-metric n=139/105/67/139 (fr slightly reduced by non-empty-reference filter).}
\label{tab:comparison}
\end{table}

\begin{table}[t]
\centering
\small
\begin{tabular}{lrcc}
\toprule
\textbf{Band} & \textbf{n} & \textbf{Cov.} & \textbf{False} \\
\midrule
high (4+ sources) & 190 & 84.7\% & 16.7\% \\
mid (3 sources)   & 146 & 87.4\% & 11.8\% \\
low (2 sources)   &  94 & 90.8\% &  5.7\% \\
rare (1 source)   &  23 & 93.5\% & 10.9\% \\
\bottomrule
\end{tabular}
\caption{Comparison by frequency band (community source coverage as proxy).}
\label{tab:bands}
\end{table}

The overall sense coverage of 87.3\% (95\% bootstrap CI: [84.7\%, 89.8\%]) with a false sense rate of 12.5\% (95\% CI: [10.1\%, 15.0\%]) shows that the model successfully preserves most reference senses while introducing few unsupported glosses (Tables~\ref{tab:comparison}--\ref{tab:bands}). Manual inspection reveals that most ``false senses'' are legitimate synonyms absent from the reference.

The main failure mode is \textbf{sense compression:} the LLM produces fewer glosses than community dictionaries (median 2 vs.\ 3), preferring the most common translation equivalent and omitting specialized or archaic senses. Manual inspection of the 57 headword--language pairs with sense coverage below 50\% reveals that sense compression (38\%) and register shift (22\%) dominate, with genuine errors (wrong meaning) accounting for only 3\%.

\paragraph{Semantic metric agreement.} Because lexical overlap can miss legitimate synonymy, Table~\ref{tab:comparison} additionally reports two model-based MT metrics on the same reference sample: BERTScore F1 with XLM-RoBERTa-large \citep{zhang2020bertscore} and COMET-DA (\texttt{wmt22-comet-da}, \citealp{rei2022comet}). Macro-averaged BERTScore~F1 is 0.924 and COMET-DA is 0.774; the language ordering tracks sense coverage, with Indonesian highest on all three metrics and German lowest on the two model-based scores. As with sense coverage, the LLM sees the community gloss as input context (\S\ref{sec:prompts}), so these values characterize context utilization rather than independent translation quality; independent quality across all 18 languages is addressed in \S\ref{sec:mt-baselines}.

For the twelve languages \textit{without} community dictionaries, the model relies on pivot translation. Section~\ref{sec:mt-baselines} addresses this gap.

\subsection{MT Baselines and Non-Community Evaluation}
\label{sec:mt-baselines}

We evaluated ten alternative systems on the same 275 headword--language evaluation pairs using a generic translation prompt (``Translate to [Language]: [text]. Be concise---dictionary style.''), deliberately simpler than our production prompt (\S\ref{sec:prompts}). We add a \textbf{format compliance} metric: whether output matches CEDICT-style gloss format (short slash-separated segments, $\leq$5 words, no sentence-ending punctuation).

\begin{table}[t]
\centering
\small
\begin{tabular}{lrccc}
\toprule
\textbf{System} & \textbf{n} & \textbf{Cov.} & \textbf{False} & \textbf{Fmt.} \\
\midrule
\textbf{Our pipeline} & \textbf{453} & \textbf{87.3} & \textbf{12.5} & \textbf{$\sim$100} \\
MiniMax (generic) & 275 & 73.0 & 27.8 & 91.0 \\
\addlinespace
Claude Sonnet 4 & 275 & 75.4 & 28.7 & 74.9 \\
Google Translate & 275 & 70.9 & 28.9 & 72.0 \\
GPT-4o-mini     & 275 & 70.4 & 32.3 & 74.2 \\
Azure Translator & 275 & 70.4 & 29.6 & 69.5 \\
DeepL            & 275 & 67.5 & 33.2 & 70.2 \\
\addlinespace
DeepSeek V3.1   & 275 & 72.0 & 36.5 & 82.1 \\
Kimi K2         & 275 & 68.0 & 35.1 & 82.1 \\
Llama 4 Scout   & 275 & 68.9 & 49.2 & 73.3 \\
Qwen3-235B      & 275 & 65.2 & 43.4 & 78.8 \\
\bottomrule
\end{tabular}
\caption{MT baseline comparison (\%). Generic prompt for all baselines. Free LLMs via Groq, Cerebras, SambaNova.}
\label{tab:mt-baseline}
\end{table}

\label{sec:backtrans}
Table~\ref{tab:mt-baseline} shows our pipeline outperforms all baselines on sense coverage (+12--22pp) and false sense rate ($-$15--37pp). The \emph{same model} (MiniMax~M2.5) with a generic prompt scores only 73.0\% coverage---14 points below our pipeline---isolating the value of dictionary-specific prompt engineering. Claude Sonnet~4, the strongest baseline at 75.4\%, still falls 12 points short. Commercial MT and frontier LLMs produce CEDICT-compatible format only 70--75\% of the time; multiple LLMs cluster in the 65--78\% range regardless of model size, suggesting pipeline design matters more than model choice.

\paragraph{Model-based metrics on the MT baselines.} We additionally score the four API-accessible baselines with BERTScore F1, COMET-DA, and reference-free CometKiwi (Chinese source) on the same 275 pairs (Table~\ref{tab:mt-baseline-semantic}). The pipeline lead extends to all three metrics: +2.1~pt BSc F1, +5.8~pt COMET-DA, +1.6~pt CometKiwi over the next-best system. Claude Sonnet~4.5 is the lowest baseline on CometKiwi (0.501) despite the highest lexical sense coverage---fluent output can correlate less tightly with the Chinese source than commercial MT trained for translation fidelity.

\begin{table}[t]
\centering
\small
\begin{tabular}{lrccc}
\toprule
\textbf{System} & \textbf{n} & \textbf{BSc.\ F1} & \textbf{COMET} & \textbf{CK} \\
\midrule
\textbf{Our pipeline} & \textbf{311} & \textbf{0.921} & \textbf{0.773} & \textbf{0.544} \\
\addlinespace
Google Translate  & 275 & 0.900 & 0.715 & 0.524 \\
Claude Sonnet 4.5 & 275 & 0.900 & 0.711 & 0.501 \\
DeepL             & 275 & 0.898 & 0.712 & 0.528 \\
GPT-4o-mini       & 275 & 0.897 & 0.702 & 0.507 \\
\bottomrule
\end{tabular}
\caption{Model-based metrics, macro-averaged over de/fr/id (same 275 pairs as Table~\ref{tab:mt-baseline}). BSc.\ F1 = BERTScore F1 (XLM-R-large); COMET = \texttt{wmt22-comet-da}; CK = reference-free \texttt{wmt22-cometkiwi-da} with the Chinese headword as source. Our-pipeline row: 3-lang macro over the same de/fr/id subset of Table~\ref{tab:comparison}. Azure and the free-LLM systems are omitted from this revision.}
\label{tab:mt-baseline-semantic}
\end{table}

Back-translation evaluation (n=1{,}571 across 13 non-community languages) yields 68.9\% sense coverage overall, with European languages strongest (Spanish 79.0\%, Portuguese 75.2\%) and Korean weakest (60.0\%). Four LLM judges (Claude Sonnet~4, Kimi~K2, GPT-OSS-120B, DeepSeek~V3.1) rate MiniMax definitions across 100 entries: accuracy \textbf{4.34}/5, naturalness \textbf{4.31}/5, completeness \textbf{3.38}/5 (mean pairwise disagreement: 0.64 points). The lower completeness aligns with sense compression (\S\ref{sec:community}). Tagalog (3.90) and Hindi (4.04) show weakest accuracy. LLM judges may share biases with the evaluated model, but provide a quality signal where none previously existed.

\paragraph{Reference-free quality across all 18 languages.} We score our pipeline output with CometKiwi (\texttt{wmt22-cometkiwi-da}, \citealp{rei2022cometkiwi}) on a 200-headword stratified sample per language (n=3{,}600). Macro-averaged CometKiwi is 0.547; the four community-reference languages average 0.562 vs.\ 0.542 for the twelve non-community languages---a 0.02 gap indicating pipeline quality does not degrade materially for non-community-anchor languages (Tagalog lowest 0.471, English highest 0.614; per-language values in Appendix~\ref{sec:appendix-cometkiwi}).

\paragraph{English-gloss-context ablation.} We compare two prompts to Llama-3.3-70B (open-weight, via Groq): \textit{direct} (Chinese-only) and \textit{pivot-context} (adds the CC-CEDICT English gloss alongside the Chinese source). On symmetric Vietnamese (n=324) and Thai (n=324) samples of headwords with CC-CEDICT glosses, the direct-minus-pivot CometKiwi delta is $-0.006$ for Vietnamese (95\% paired-bootstrap CI $[-0.020, +0.009]$, B=2000) and $-0.024$ for Thai (CI $[-0.039, -0.010]$). Vietnamese straddles zero (no detectable effect); Thai is entirely negative, indicating pivot context measurably helps Thai by $\sim$0.024 CometKiwi points ($\sim$5\% of the direct score). Symmetric gloss agreement (mean of forward+backward sense coverage between conditions) is 61.3\% (vi) and 40.8\% (th), ruling out the delta being noise between near-identical outputs. Two caveats: (i)~CometKiwi is English-heavy in training so absolute low-resource scores carry uncertainty, though within-language contrast largely cancels this bias; and (ii)~this is pivot-as-context (source retained, gloss added), not the classical A$\to$MT(en)$\to$MT(B) cascade to which prior negative pivot-MT results apply.

\paragraph{Per-language error taxonomy.} A full-corpus analysis (Table~\ref{tab:error-taxonomy}, Appendix~\ref{sec:error-taxonomy}) quantifies mono-gloss rate (55\% en $\rightarrow$ 74.6\% id/tl), CJK-character rate ($\sim$1\% non-CJK targets, 87\% ja, 2.7\% ko), and Chinese-headword echo (13\% for ja).

\paragraph{Second-model verification.} Directly answering the reviewer suggestion to have a second model verify the no-community-reference backfill, we asked an independent open-weight model (Gemma-4-31B via Cerebras) to rate each MiniMax gloss 1--5 for semantic preservation of the Chinese source. Across 14 non-community-reference target languages (n=2{,}800 rated, 200 per lang stratified by frequency band), the macro mean rating is \textbf{4.47/5} and only \textbf{9.6\%} of glosses are flagged as low-quality (rating $<$3). Per-language means range from 4.14 (Tagalog) to 4.65 (Portuguese/Russian), with every language above 4.0. The two lowest per-language means (Tagalog 4.14, Thai 4.35) coincide with the lowest CometKiwi scores in Table~\ref{tab:cometkiwi} (0.471, 0.521)---two independent open-weight quality-estimation methods agree on which languages need the most improvement.

\subsection{Source Priority and Coverage}

For headwords with definitions from multiple sources, we apply a priority hierarchy: community dictionaries (CC-CEDICT, HanDeDict, CFDICT, CC-CIDICT, CHEDICC, JMdict) $>$ Wiktextract $>$ MiniMax; human-curated definitions take precedence. Every one of the 428,073 headwords has definitions in all 18 target languages---7{,}705{,}314 slots total, 100\% coverage; 7{,}697{,}901 slots carry an LLM-generated gloss and 738{,}777 additionally carry a community-dictionary gloss. The dictionary is further augmented with 511{,}514 Cantonese and Hokkien dialect forms. Total LLM inference cost: {\raise.17ex\hbox{$\scriptstyle\sim$}}\$146 at published pay-as-you-go rates.


%% ============================================================
\section*{Limitations}
\label{sec:limitations}

\textbf{Glosses, not definitions.} The resource produces translation equivalents, not lexicographic definitions with usage notes or grammatical patterns; polysemous entries carry flat gloss lists (no sense discrimination), and inflecting languages lack grammatical metadata (German gender, Russian aspect).

\textbf{Evaluation limitations.} Sense coverage is an automated lexical-overlap proxy, and partially circular for languages where the community gloss appears in prompt context (\S\ref{sec:community}). The model-based metrics (BERTScore, COMET-DA, CometKiwi) mitigate lexical-overlap surface bias but inherit the biases of their multilingual encoders; CometKiwi is less well-calibrated for low-resource targets. LLM-as-judge scores (\S\ref{sec:backtrans}) may share biases with the evaluated model. Human evaluation across all 18 languages remains future work.

\textbf{Pivot bias and single-model homogeneity.} All translations are mediated through English context, inheriting English sense boundaries; all 7.7M glosses originate from MiniMax~M2.5, so model-specific biases propagate uniformly. Reproducibility is nevertheless supported by M2.5's open weights (hosted on Fireworks, NVIDIA~NIM; self-hostable).

\textbf{Scope.} The 428K headwords are inherited without curation; the 18 target languages cover major families but omit Malay, Bengali, and Swahili.


%% ============================================================
\section{Future Work}
\label{sec:future}

\textbf{From glossary to dictionary.} Sense-discriminated entries with usage guidance and grammatical enrichment (gender, aspect, POS).

\textbf{Computational loanword detection.} The dictionary's Hokkien source forms and Southeast Asian targets enable a machine-readable atlas of Chinese vocabulary diffusion---linking \textit{tāu-hū} to Indonesian \textit{tahu}, Tagalog \textit{taho}, and Thai \foreignlanguage{thai}{เต้าหู้}---across maritime Southeast Asia \citep{chanyap1980}.

\textbf{Under-resourced languages.} The $\sim$\$146 cost makes the pipeline relevant for languages without institutional lexicography support (Zhuang, Tibetan, Uyghur, Yi within China alone). Applying it inherits pivot-translation limitations \citep{haddow2022lowresource}; the ablation in \S\ref{sec:mt-baselines} shows this risk is language-dependent.

\textbf{Community-in-the-loop evaluation.} Native-speaker evaluation across all 18 targets---particularly by diaspora communities---would establish gold-standard metrics. A realistic instrument (2 annotators, 3--5 languages, $\sim$100 stratified entries; accuracy/naturalness/completeness on 5-point scales) is deferred pending recruitment funding; the model-based metrics in \S\ref{sec:community}--\S\ref{sec:mt-baselines} are the strongest automated proxy.


%% ============================================================
\section{Conclusion}
\label{sec:conclusion}

A single open-weight LLM (MiniMax M2.5) can produce a multilingual Chinese glossary of 428,073 headwords $\times$ 18 languages for \${\raise.17ex\hbox{$\scriptstyle\sim$}}146. Enablers: batched cross-language prompting, language-specific rules that reduce script contamination to $\sim$0.83\%, a deterministic script-validator gate, and context-aware backfill (100\% coverage). Combined with 511K Cantonese/Hokkien dialect forms, this is the first open multilingual Chinese lexical resource covering the major languages of Europe, East Asia, Southeast Asia, and the Middle East. The dictionary is deployed in production in a Chinese learning application.

%% ============================================================
\section*{Data and Code Availability}
\label{sec:availability}
The full \textit{Dictionarium Sinicum} SQLite database (2\,GB, 428{,}073 headwords $\times$ 18 languages, SHA-256 \texttt{e647e163\ldots8ea512}) is archived at Harvard Dataverse under CC BY-SA 4.0: \url{https://doi.org/10.7910/DVN/HJWMKI}. For the anonymous review period, an anonymized preview URL (stripping author metadata) is available at \url{https://dataverse.harvard.edu/previewurl.xhtml?token=a53578a1-5c98-483e-9a86-49b653d8c5a3}. All pipeline code, evaluation scripts, and reproducibility artifacts (Apache-2.0) are mirrored anonymously at \url{https://anonymous.4open.science/r/arr26paperbmirror/} for the review period; the production repository URL will be provided in the camera-ready.


%% ============================================================
% Acknowledgments: omit for review, include in camera-ready
%\section*{Acknowledgments}
%Claude (Anthropic) was used extensively as a development partner throughout this project: writing and debugging pipeline code, iterating on prompt engineering, performing data analysis, and assisting with manuscript preparation. Dictionary definitions were generated by MiniMax M2.5 as described in Section~4. The author is solely responsible for all research decisions, claims, and errors.

%% ============================================================
\bibliography{references}

%% ============================================================
\appendix

\section{Prompt Templates}
\label{sec:appendix-prompts}

\subsection{System Prompt (production)}

\begin{lstlisting}
You are a professional multilingual Chinese
lexicographer producing dictionary-style
definitions in 18 languages.

Rules:
- Output EXACTLY one line per language in
  format "xx: def1/def2"
- Be concise: dictionary style, not full
  sentences
- Maximum 5 glosses per entry
- CRITICAL: Every non-Chinese definition must
  contain ZERO Chinese characters.

Language-specific rules:
- ja: Provide the MEANING in Japanese, not
  just kanji echo or kana reading.
- ko: Write in Hangul only. Do NOT mix in
  Chinese characters.
- vi: Write in Vietnamese with proper
  diacritics. No Chinese characters.
- ar: Write in Arabic script only. No Latin.
- th: Write in Thai script with tone marks.
- hi: Write in Devanagari script.
- fa: Write in Persian script only.
- nl/pt/it: Write in standard Dutch/
  Portuguese/Italian.
\end{lstlisting}

\subsection{User Prompt (with corpus examples)}

\begin{lstlisting}
Chinese: 銀行 / 银行
Pinyin: yin hang
POS: noun

Existing definitions:
  English: bank/financial institution
  German: Bank/Geldinstitut
  French: banque/etablissement bancaire

Example sentences:
  中国人民银行决定下调金融机构贷款和存款基准利率。
  这家银行在全国设有两千多个分支机构。

Produce definitions for each language below.
Output EXACTLY one line per language.

en:
de:
fr:
es:
sv:
ja:
ko:
ru:
id:
vi:
tl:
fa:
nl:
pt:
ar:
th:
hi:
it:
\end{lstlisting}

\subsection{Backfill Prompt (context-aware)}

\begin{lstlisting}
Fill in ONLY the missing language definitions.
Existing translations are provided for
consistency -- match their style and specificity.
Do NOT repeat or modify existing translations.

SCRIPT PURITY (strictly enforced -- output
that violates these rules will be rejected):
Each definition must use ONLY the script of
its target language.

1. 三角褲 / 三角裤
   Pinyin: san jiao ku
   Existing translations:
     en: briefs/underwear
     de: Slip/Unterhose
     fr: slip/calecon
     [... 14 more languages ...]
   MISSING -- fill these:
     vi:
\end{lstlisting}

\section{Sample Definitions}
\label{sec:appendix-samples}

\subsection{Idiom: 画蛇添足}
\textit{huà shé tiān zú} ``draw legs on a snake''

\begin{lstlisting}
en: to improve something unnecessarily;
    to gild the lily
de: der Schlange Fuesse hinzumalen;
    etwas Ueberfluessiges tun
fr: dessiner un serpent et lui ajouter
    des pattes; faire du zele
es: dibujar una serpiente y anadirle
    patas; hacer algo innecesario
ko: 뱀에게 발을 그리다;
    불필요한 것을 추가하다
vi: ve ran them chan;
    them thu khong can thiet
\end{lstlisting}

\subsection{Cultural term: 春节}
\textit{chūn jié} ``Spring Festival''

\begin{lstlisting}
en: Spring Festival (Chinese New Year)
de: Chinesisches Neujahrsfest
fr: Nouvel An chinois/Fete du printemps
id: Festival Musim Semi
ja: 春節(中国の旧正月)
ko: 춘제/설날
vi: Tet Nguyen dan
\end{lstlisting}

\subsection{Loanword chain: 豆腐}
\textit{dòufu} ``tofu''

\begin{lstlisting}
en: tofu / bean curd
de: Tofu / Bohnenkäse
id: tahu          <- Hokkien tau-hu
tl: tokwa         <- Hokkien tau-koan
vi: dau phu       <- Sino-Vietnamese
ko: 두부          <- Sino-Korean
ja: とうふ        <- Sino-Japanese
\end{lstlisting}

This entry illustrates how the Hokkien maritime trade vocabulary seeded cognates across Southeast Asian languages: \textit{tāu-hū} (Hokkien) → \textit{tahu} (Indonesian/Malay), \textit{taho/tokwa} (Tagalog), \textit{đậu phụ} (Vietnamese, Sino-Vietnamese route).

%% ============================================================
\section{Reference-Free Quality Estimates}
\label{sec:appendix-cometkiwi}

Table~\ref{tab:cometkiwi} reports per-language CometKiwi scores (\S\ref{sec:mt-baselines}) on the 200-headword stratified sample per language.

\begin{table}[h]
\centering
\small
\begin{tabular}{lc|lc}
\toprule
\textbf{Lang} & \textbf{CK} & \textbf{Lang} & \textbf{CK} \\
\midrule
en$^{\dagger}$ & 0.614 & id$^{\dagger}$ & 0.543 \\
de$^{\dagger}$ & 0.534 & ja & 0.546 \\
fr$^{\dagger}$ & 0.555 & ko & 0.557 \\
es & 0.558 & ru & 0.568 \\
it & 0.570 & vi & 0.554 \\
pt & 0.554 & tl & 0.471 \\
nl & 0.540 & fa & 0.535 \\
sv & 0.540 & th & 0.521 \\
      &       & ar & 0.503 \\
      &       & hi & 0.589 \\
\midrule
\multicolumn{4}{c}{\textbf{Macro: 0.547} (ref 0.562 / non-community 0.542)} \\
\bottomrule
\end{tabular}
\caption{Reference-free CometKiwi scores (higher is better) on 200 stratified headwords per language (n=3{,}600 total). $\dagger$ = language with a community reference dictionary. Left column groups European targets; right column groups East Asian, Southeast Asian, and Middle Eastern targets.}
\label{tab:cometkiwi}
\end{table}

\section{Per-Language Error Taxonomy}
\label{sec:error-taxonomy}

Table~\ref{tab:error-taxonomy} reports automated per-language error signals across the full 428,073-headword MiniMax output.

\begin{table}[h]
\centering
\small
\begin{tabular}{lrrrrr}
\toprule
\textbf{Lang} & \textbf{Mono} & \textbf{Mean} & \textbf{Med.} & \textbf{CJK} & \textbf{Zh-}\\
              & \textbf{\%}   & \textbf{seg.} & \textbf{len}  & \textbf{\%}  & \textbf{echo\%}\\
\midrule
en  & 55.4 & 1.56 & 25 & 1.07 & 0.57 \\
de  & 67.4 & 1.40 & 23 & 1.02 & 0.51 \\
fr  & 69.6 & 1.37 & 22 & 1.05 & 0.53 \\
id  & 74.6 & 1.31 & 20 & 1.11 & 0.57 \\
es  & 70.3 & 1.36 & 22 & 1.06 & 0.54 \\
sv  & 71.5 & 1.35 & 18 & 1.02 & 0.53 \\
ja  & 61.7 & 1.45 &  9 & 87.45 & 13.01 \\
ko  & 68.0 & 1.40 &  7 & 2.73 & 0.53 \\
ru  & 71.0 & 1.35 & 21 & 0.98 & 0.49 \\
vi  & 71.6 & 1.34 & 18 & 0.99 & 0.53 \\
tl  & 73.5 & 1.32 & 22 & 1.04 & 0.55 \\
fa  & 71.3 & 1.35 & 16 & 0.72 & 0.35 \\
nl  & 68.4 & 1.38 & 20 & 1.03 & 0.52 \\
pt  & 69.1 & 1.37 & 21 & 1.07 & 0.52 \\
ar  & 65.1 & 1.41 & 17 & 1.46 & 0.40 \\
th  & 68.3 & 1.38 & 18 & 0.80 & 0.36 \\
hi  & 67.1 & 1.39 & 18 & 0.77 & 0.39 \\
it  & 68.8 & 1.37 & 22 & 1.02 & 0.52 \\
\bottomrule
\end{tabular}
\caption{Per-language error signals from the full 428,073-headword MiniMax output. Mono~\% = fraction of definitions with only one slash-separated gloss (sense-compression proxy). Mean~seg. = mean glosses per definition. Med.~len = median character length. CJK~\% = fraction containing any CJK character (script contamination for non-CJK targets; expected for ja/ko). Zh-echo~\% = fraction containing the Chinese headword string (metalinguistic reference, common in ja).}
\label{tab:error-taxonomy}
\end{table}

%% ============================================================
\section{Response to March 2026 Reviewers}
\label{sec:response-letter}

\textbf{Previous review round:} ARR March 2026, submission ID \texttt{sVPHR4PSzR}, reviewers \texttt{gKVr}, \texttt{SamJ}, \texttt{cky5}, meta-review \texttt{zfRe}.

We are grateful to the March 2026 reviewers for their detailed feedback. This appendix maps every substantive objection from the previous round to the specific manuscript changes that address it, and preempts likely questions from this round by anticipating them below.

\subsection{Changes since ARR March 2026}
\label{sec:response-changes}

\paragraph{SamJ: ``Consider including COMET or BERTScore alongside lexical overlap to better support claims and capture semantic accuracy.''}

\textbf{Addressed.} \S\ref{sec:community} (Table~\ref{tab:comparison}) now reports BERTScore~F1 and COMET-DA alongside sense-coverage and false-sense rate on the four community reference languages (en, de, fr, id; macro F1 0.924, COMET-DA 0.774). \S\ref{sec:mt-baselines} (Table~\ref{tab:cometkiwi}) reports reference-free CometKiwi across all 18 target languages with the Chinese headword as the anchor (macro 0.547; ref langs 0.562, non-community 0.542). \S\ref{sec:mt-baselines} (Table~\ref{tab:mt-baseline-semantic}) reports the same three model-based metrics for four API-accessible MT baselines (Claude Sonnet~4.5, GPT-4o-mini, Google Translate, DeepL); the pipeline leads on all three metrics.

\paragraph{SamJ: ``For the backfill phase, did the authors consider a second-model verification approach?''}

\textbf{Addressed.} \S\ref{sec:mt-baselines} ``Second-model verification'' paragraph reports an independent Gemma-4-31B (open-weight, Cerebras) rating of MiniMax outputs across all 14 non-community-reference languages, 200 stratified headwords each (n=2{,}800 total). Macro mean rating 4.47/5; only 9.6\% of glosses are flagged as low-quality (rating $<$3). The two lowest per-language means (Tagalog 4.14, Thai 4.35) coincide with the lowest CometKiwi scores (0.471, 0.521)---two independent open-weight quality-estimation methods agree on which languages need the most improvement.

\paragraph{SamJ: PDF hyphenation typos.}

\textbf{Addressed.} The hyphenation glitches were fixed in a dedicated Tier~1 pass: \texttt{microtype} factor~1050, explicit \texttt{\textbackslash lefthyphenmin=3 \textbackslash righthyphenmin=3}, targeted \texttt{\textbackslash hyphenation\{\}} list of $\sim$50 words. Over 20 two-letter fragment splits in the March~2026 version are now zero.

\paragraph{gKVr / cky5: ``Software: 1'' --- no public code.}

\textbf{Addressed.} The pipeline and evaluation code are released at an anonymized mirror (see the Data and Code Availability section). The production repository is Apache-2.0 licensed and includes: the full three-phase translation pipeline, all six eval scripts (community comparison, MT baselines, semantic metrics, pivot ablation, error taxonomy, second-model verification), the dictmaster SQLite schema and importers, and 189~passing unit tests.

\paragraph{gKVr: ``Single proprietary model.''}

\textbf{Addressed as a paper-clarity issue.} MiniMax M2.5 is open-weight; it is hosted on multiple third-party inference providers (Fireworks, NVIDIA NIM) and can be self-hosted. The Abstract, \S\ref{sec:intro}, \S\ref{sec:pipeline}, the Limitations section, and \S\ref{sec:future} now explicitly call this out. Additionally, \S\ref{sec:mt-baselines} (English-gloss-context ablation) uses Llama-3.3-70B (open-weight, Meta) to demonstrate the pipeline pattern works with a second open-weight model family.

\paragraph{gKVr: ``Limited eval scope on the 12 non-community-reference languages.''}

\textbf{Addressed by two independent evaluations covering all 18 languages} (not just the reference-language subset): (a) reference-free CometKiwi with Chinese as source (\S\ref{sec:mt-baselines} Table~\ref{tab:cometkiwi}) shows the non-community-language macro (0.542) is only 0.02 below the reference-language macro (0.562); (b) second-model verification (\S\ref{sec:mt-baselines}, 14 non-community-ref languages including Spanish and Japanese) gives macro 4.47/5 with all 14 languages above 4.0.

\paragraph{gKVr: ``Anglocentric bias / pivot criticism.''}

\textbf{Addressed by explicit pivot-vs-direct ablation} (\S\ref{sec:mt-baselines} English-gloss-context ablation paragraph). Vietnamese and Thai, same open-weight model (Llama-3.3-70B), with vs.\ without CC-CEDICT English gloss as prompt context. At properly powered sample sizes (n=324 per language) with percentile-bootstrap 95\% CIs, the direct-minus-pivot CometKiwi delta is $-0.006$ for Vietnamese (CI $[-0.020, +0.009]$) and $-0.024$ for Thai (CI $[-0.039, -0.010]$). The Vietnamese interval straddles zero (no detectable effect); the Thai interval is entirely negative, indicating pivot context measurably helps for Thai. Symmetric gloss agreement is 61.3\%~(vi) and 40.8\%~(th), ruling out the delta being noise between near-identical outputs. Finding: pivot effect is language-dependent, not a uniform quality loss.

\paragraph{AC: ``Error analysis / hallucination analysis / per-language breakdown.''}

\textbf{Addressed.} \S\ref{sec:mt-baselines} ``Per-language error taxonomy'' paragraph points to Appendix~\ref{sec:error-taxonomy}, which reports full-corpus automated per-language error signals (mono-gloss rate, mean segments, median length, CJK-character inclusion, Chinese-headword echo) for all 18 languages. Notable patterns: mono-gloss rate ranges 55\% (English) to 74.6\% (Indonesian/Tagalog); CJK inclusion is $\sim$1\% for non-CJK targets except Japanese (87\%, legitimate kanji) and Korean (2.7\%, low-level hanja mixing).

\paragraph{cky5: Human evaluation.}

\textbf{Framed as future work pending funding} (\S\ref{sec:future}). A realistic native-speaker instrument (2 annotators, 3--5 languages, $\sim$100 stratified entries, 5-point accuracy/naturalness/completeness scales) would raise soundness scores substantially. The model-based metrics in \S\ref{sec:community}--\S\ref{sec:mt-baselines} are the strongest automated proxy currently available; we prioritized deploying them for this revision rather than the more expensive human study.

\subsection{Anticipated questions this round}
\label{sec:response-anticipated}

We ran an unofficial pre-review round with three independent large language models (Claude Fable~5, xAI Grok, Google Gemini) prior to submission. Their strongest recurring critiques and our responses:

\paragraph{Q1: Isn't the +12--22pp sense-coverage lead over MT baselines confounded because the pipeline sees community glosses in prompt context but the baselines don't?}

Yes, and the manuscript acknowledges this. The Abstract, \S\ref{sec:community} caveat paragraph, and the framing around Tables~\ref{tab:mt-baseline}/\ref{tab:mt-baseline-semantic} have been softened for this revision: the +12--22pp gap is now explicitly attributed partly to the pipeline's access to community-reference glosses in its input context, not to independent generation capability alone. The paper distinguishes two claims:

\begin{itemize}
\item The community-reference languages (en/de/fr/id) show \textbf{context utilization}: how faithfully does the model preserve reference senses when given them? (Not a fair test of independent quality.)
\item The 14 non-community-reference languages show \textbf{pivot translation quality} under a genuine open-ended MT setting; these are scored with reference-free CometKiwi and independent Gemma-4-31B verification.
\end{itemize}

The +12--22pp headline claim was retained because the \emph{same model} (MiniMax M2.5) with a generic zero-context prompt scores 14 points lower than our pipeline on the same 275-pair evaluation, isolating a real prompt-engineering effect.

\paragraph{Q2: The pivot ablation uses only $\sim$33 pairs per language after CC-CEDICT filtering. Isn't that statistically underpowered?}

\textbf{We agree, and rerunning changes the finding.} The original n=33 suggested pivot helped Vietnamese by $+0.067$ CometKiwi and was neutral for Thai. Rerun at n=324 (vi) and n=324 (th) with percentile-bootstrap 95\% CIs shows the delta shrinks by an order of magnitude AND the Vietnamese CI straddles zero: Vietnamese direct-minus-pivot $\Delta = -0.006$ (CI $[-0.020, +0.009]$), Thai $\Delta = -0.024$ (CI $[-0.039, -0.010]$). The earlier ``pivot helps Vietnamese'' finding was noise, not signal. The revised \S\ref{sec:mt-baselines} paragraph reports the honest finding: at properly powered sample sizes, English pivot has no statistically detectable effect for Vietnamese and a small measurable positive effect for Thai. Notably, the gloss agreement is only 61\%~(vi) / 41\%~(th), meaning pivot and direct produce lexically different translations of comparable reference-free quality --- the model uses the pivot context but the target-language sense boundaries land in similar semantic space either way.

\paragraph{Q3: What is the sense-coverage performance on a held-out set of headwords whose community glosses were NEVER supplied in the prompt?}

This is a fair follow-up. We defer a true held-out evaluation to a future revision because the production pipeline unconditionally supplies community glosses for all reference-language slots (this is design intent --- context utilization is the point). A retrospective held-out slice would require rerunning the pipeline on a masked subset and would fully replicate the $\sim$\$146 cost per run. We commit to running one such slice for the camera-ready if the paper is accepted.

\paragraph{Q4: Why MiniMax M2.5 over other open-weight models that also passed the 5-headword smoke test (Kimi K2, Nemotron)?}

Table~\ref{tab:models} documents the selection criteria: format compliance, script contamination, throughput, and cost. MiniMax M2.5 was the only model that produced zero issues on the 5-headword suite across all four criteria. Kimi K2 and Nemotron both passed format compliance but had non-zero contamination rates on the smoke test. The paper adds an explicit note that a full-scale rerun with a different generator (e.g., Qwen3-235B or Llama-3.3-70B) is future work; the Llama-3.3-70B pivot ablation and the Gemma-4-31B verification pass demonstrate that the pipeline pattern generalizes across open-weight model families.

\paragraph{Q5: Are there residual data-integrity issues after the three backfill passes?}

No. A post-cleanup audit (\S\ref{sec:quality}) shows every headword $\times$ language slot is filled, and the deterministic script validator has an audit tab in the released code. The paper's earlier claim of ``90.6\% single-character coverage'' was a snapshot from an intermediate run; subsequent backfill iterations with alternative prompting for metalinguistic definitions closed those gaps. The current DB (released at Harvard Dataverse, doi:10.7910/DVN/HJWMKI) shows 100.0\% coverage across all 7{,}709{,}256 slots.

\paragraph{Q6: Missing prior art?}

Gemini flagged three references worth adding for context and completeness: Bond \& Paik (2012) on wordnet licensing, Haddow et al.\ (2022) on low-resource MT, and Wang et al.\ (2023) on document-level MT with LLMs. All three are now cited (\S\ref{sec:related}, \S\ref{sec:future}).

\subsection{Known remaining limitations}
\label{sec:response-limitations}

Documented in the Limitations section. Highest-priority items for future work: (a) native-speaker human evaluation on 3--5 non-community-ref languages (deferred pending recruitment funding); (b) full-pipeline rerun with an alternative open-weight generator (deferred pending $\sim$\$150 compute budget); (c) sense discrimination and grammatical enrichment (multi-year research program).

\subsection{Reproducibility}
\label{sec:response-reproducibility}

\begin{itemize}
\item Anonymized code mirror at \texttt{anonymous.4open.science} (link in \S\ref{sec:availability}).
\item Full 2\,GB SQLite database archived at Harvard Dataverse under CC~BY-SA~4.0 (\texttt{doi:10.7910/DVN/HJWMKI}; see \S\ref{sec:availability}).
\item 189/189 unit tests pass on the released repository.
\item All eval scripts accept \texttt{--seed} for reproducible sampling; every reported metric can be re-derived by running the corresponding script on the released database.
\end{itemize}

\end{document}
